\chapter{Metodologia assistita da agenti AI}
\label{cap:metodologia-ai}

\intro{Il presente capitolo descrive la metodologia di lavoro adottata durante lo stage, incentrata sull'uso di un agente di codifica basato su un modello linguistico di grandi dimensioni (\gls{llm}) come \emph{pair engineer} strutturato. Vengono illustrati gli ambiti di utilizzo, il workflow \emph{Plan $\rightarrow$ Confirm $\rightarrow$ Implement}, i benefici in termini di tempo e le mitigazioni adottate per i limiti noti.}\\

\section{Contesto e motivazione}

La pianificazione dello stage prevedeva una durata di 300 ore distribuite su 8 settimane,
a fronte di un ambito ampio: apprendimento di tecnologie nuove (OpenTelemetry SDK .NET,
stack \emph{Grafana} LGTM, PromQL/TraceQL/LogQL, Serilog OTLP, \emph{Spec-Kit}), realizzazione
di prototipi, instrumentazione di più moduli del framework \emph{Sia.Web} e produzione della
relativa documentazione. Il rapporto tra l'ampiezza del tema e il tempo disponibile ha
suggerito l'adozione di un agente AI non come generatore one-shot di codice, ma come
\emph{pair engineer} affiancato allo sviluppatore, capace di accelerare le attività ripetitive
e di replicare pattern già validati su nuovi moduli.

La scelta è ricaduta su \emph{Claude Code}, l'agente di codifica di Anthropic integrato nel
terminale e in grado di interagire con il \emph{filesystem}, eseguire comandi e invocare
strumenti esterni tramite \gls{mcp}.

\section{Ambiti di utilizzo}

L'agente è stato utilizzato in tre ambiti distinti, corrispondenti a tre momenti differenti
dello stage.

\subsection{Studio delle tecnologie}

Nelle prime tre settimane, dedicate allo studio delle tecnologie, l'agente è stato impiegato
per sintetizzare la documentazione ufficiale di \gls{otel}, Tempo, Loki, Prometheus, e per
produrre confronti mirati tra le convenzioni \emph{naming} di \gls{otel} e quelle esposte dagli
\emph{exporter} Prometheus (ad esempio il trattamento delle \texttt{unit} e l'espansione del
suffisso \texttt{\_milliseconds}/\texttt{\_seconds}). Tale attività ha permesso di ridurre in
modo sensibile il tempo stimato di \emph{onboarding}, concentrando lo studio sui soli punti
critici per il progetto invece che sulla documentazione integrale di ciascuna tecnologia.

\subsection{Prove di implementazione}

Durante la fase di progettazione, l'agente è stato usato per produrre rapidamente varianti di
\emph{proof of concept} (\gls{poc}) su scelte ancora da validare: \emph{sampler} \emph{head-based}
vs \emph{tail-based}, esportazione diretta dall'applicazione vs passaggio tramite
\emph{Collector}, esempi di istrumentazione automatica vs istrumentazione manuale.
Avere a disposizione codice compilabile in pochi minuti ha consentito di \emph{provare e
scartare} alternative, invece di decidere a priori sulla base della sola documentazione.

\subsection{Propagazione della telemetria agli altri moduli}

L'ambito in cui l'uso dell'agente si è rivelato più produttivo è stato la propagazione
dell'instrumentazione ai moduli del framework (\emph{Sia.Base}, \emph{Sia.Auth},
\emph{Sia.Grid}, \emph{Sia.Crud}, \ldots). Per evitare che la replicazione del pattern di
instrumentazione su moduli diversi introducesse incoerenze (naming, cardinalità, attributi
di \emph{resource}), sono state definite \emph{Claude skill} dedicate:

\begin{itemize}
    \item \texttt{sia-telemetry-setup}: configurazione del pipeline OTel, \emph{sampler},
        \emph{exporter}, \emph{appsettings};
    \item \texttt{sia-telemetry-traces}: \emph{span}, \emph{tag}, \emph{helper} \texttt{TrackQuery};
    \item \texttt{sia-telemetry-metrics}: metriche, \emph{naming}, cardinalità, pattern
        \texttt{RecordMetricError};
    \item \texttt{sia-telemetry-logs}: Serilog OTLP, correlazione \emph{log}-\emph{trace};
    \item \texttt{sia-telemetry-grafana}: \emph{dashboard} Grafana, query PromQL/TraceQL/LogQL,
        strumenti \gls{mcp}.
\end{itemize}

Ogni \emph{skill} è un contratto procedurale, non un semplice \emph{prompt}, che descrive
convenzioni, \emph{anti-pattern} e \emph{template} di codice per uno specifico aspetto. L'agente
carica la \emph{skill} appropriata quando riceve un task pertinente, garantendo che
moduli instrumentati in momenti diversi rispettino le stesse convenzioni.

\section{Workflow \emph{Plan $\rightarrow$ Confirm $\rightarrow$ Implement}}
\label{sec:workflow-plan}

Per evitare i rischi tipici della generazione automatica di codice (\emph{sprawl},
allucinazioni, \emph{overfitting} al primo caso visto), è stato adottato un \emph{workflow}
strutturato in tre fasi.

\begin{enumerate}
    \item \textbf{Plan.} Per ogni task non banale, l'agente produce un file
        \texttt{plan.md} nella cartella\\
        \texttt{.claude/implementations/\{yy\}/\{MM\}/\{dd\}\_\{titolo\}/} che descrive
        obiettivo, componenti coinvolti, modifiche pianificate file per file, assunzioni e
        domande aperte. Nessun codice viene prodotto in questa fase.
    \item \textbf{Confirm.} Lo sviluppatore revisiona il \texttt{plan.md}, può editarlo
        direttamente, integrarlo o rifiutarlo. L'agente si ferma \emph{esplicitamente} in
        attesa di approvazione.
    \item \textbf{Implement.} Solo dopo approvazione l'agente scrive un \texttt{task.md}
        (\emph{checklist} operativa) e implementa in ordine, chiudendo ogni sezione con uno
        \emph{stato} \texttt{[x]}.
\end{enumerate}

Il \emph{workflow} produce un artefatto tracciabile per ogni intervento: \texttt{plan.md} e
\texttt{task.md} restano nel repository e fungono da registro decisionale.

---------------------- AGGIUNGERE ESEMPIO DI PLAN.MD E TASK.MD --------------

\section{Model Context Protocol}
\label{sec:mcp-methodology}

Il \gls{mcp} è il protocollo, introdotto nel 2024, che consente a un agente basato su modelli
linguistici di interagire in modo strutturato con servizi esterni, ottenendo dati dinamici e
contestuali. Nel progetto è stato utilizzato per collegare \emph{Claude Code} all'istanza
\emph{Grafana} di sviluppo.

\subsection{Obiettivi}
\begin{itemize}
    \item \textbf{Accesso a dati runtime}: consentire all'agente di interrogare dati di
        telemetria in tempo reale (metriche effettivamente presenti in Prometheus, \emph{label}
        disponibili, dashboard esistenti).
    \item \textbf{Contesto dinamico}: fornire informazioni aggiornate sul comportamento
        dell'applicazione in esecuzione.
    \item \textbf{Integrazione modulare}: permettere l'aggiunta di nuovi servizi (Tempo, Loki,
        Prometheus) senza modificare l'agente.
    \item \textbf{Standardizzazione}: definire un'interfaccia comune per la comunicazione tra
        agenti e strumenti esterni.
\end{itemize}

\subsection{Uso degli strumenti MCP}
Gli strumenti \gls{mcp} principali usati nel progetto sono:
\begin{itemize}
    \item \texttt{search\_dashboards}, \texttt{get\_dashboard\_by\_uid},
        \texttt{update\_dashboard} per creare e versionare le \emph{dashboard};
    \item \texttt{list\_prometheus\_metric\_names}, \texttt{list\_prometheus\_label\_values},
        \texttt{query\_prometheus} per validare \emph{naming} e \emph{label} delle metriche;
    \item \texttt{list\_loki\_label\_values}, \texttt{query\_loki\_logs} per costruire e
        validare query LogQL sui pannelli di log.
\end{itemize}

Dopo l'introduzione del protocollo, la qualità delle dashboard prodotte dall'agente è
migliorata in modo marcato, grazie alla possibilità di verificare la presenza effettiva di
metriche e \emph{label} prima di comporre la query. È stata comunque necessaria
un'integrazione di direttive specifiche per prevenire allucinazioni residue: uso di metriche
o \emph{label} inesistenti, applicazione di filtri su attributi presenti sulle \emph{span} ma
non sulle metriche (es.\ \texttt{componentId}), e organizzazione dei pannelli
secondo il pattern \emph{durata $\rightarrow$ throughput $\rightarrow$ errori $\rightarrow$ log}.

\section{Confronto pianificazione vs esecuzione}

L'uso dell'agente ha prodotto un beneficio misurabile soprattutto nelle fasi a forte
componente ripetitiva (studio, \gls{poc}, propagazione del pattern di instrumentazione), e un
beneficio meno marcato in quelle a forte componente decisionale (scelte architetturali,
analisi di cardinalità). Nel complesso, il tempo risparmiato nelle fasi assistite è stato
riinvestito nelle attività non automatizzabili: revisione semantica del codice prodotto,
validazione delle \emph{dashboard} contro Prometheus, affinamento iterativo delle
\emph{skill}.

\section{Limiti e mitigazioni}

L'uso di un agente \gls{llm} comporta rischi noti che sono stati mitigati in modo sistematico.

\begin{description}
    \item[Allucinazioni su metriche/label.] L'agente tende a inventare nomi di metriche o
        \emph{label} plausibili ma inesistenti. \textbf{Mitigazione}: regola obbligatoria di
        validazione tramite \gls{mcp} Grafana (\texttt{list\_prometheus\_metric\_names},
        \texttt{list\_prometheus\_label\_values}) prima di chiudere ogni task dashboard.
    \item[Overfitting del prompt.] Le \emph{skill} iniziali, scritte a partire dal primo
        modulo instrumentato (\texttt{Sia.Base}), contenevano riferimenti troppo specifici
        (nomi di metriche, \emph{tag}, stored procedure). \textbf{Mitigazione}: riscrittura
        delle \emph{skill} in termini generali, con placeholder \texttt{\{modulo\}} e
        riferimenti ai moduli canonici solo come esempi.
    \item[Perdita di contesto su sessioni lunghe.] La \emph{context window} dell'\gls{llm} è
        finita; su task lunghi, l'agente perde traccia di decisioni prese in precedenza.
        \textbf{Mitigazione}: memorizzazione persistente tramite \texttt{MEMORY.md} per
        progetto e per agente, \texttt{plan.md} e \texttt{task.md} come registro decisionale
        per ogni task.
    \item[Divergenza dei pattern tra moduli.] Sviluppatori diversi, o sessioni diverse, possono
        produrre variazioni del pattern di instrumentazione. \textbf{Mitigazione}:
        centralizzazione dei pattern obbligatori nell'agente \texttt{sia-telemetry-developer}
        (es.\ counter errori controller via \texttt{IControllerDiagnostics} con Adapter, uso
        obbligatorio della \emph{template variable} \texttt{bucket} nelle dashboard).
\end{description}

\section{Considerazioni conclusive}

L'uso di un agente AI come \emph{pair engineer} strutturato, non come strumento di
\emph{autocompletion}, ha reso fattibile in 8 settimane sia l'implementazione della
telemetria sul modulo pilota sia la sua propagazione trasversale a più moduli del framework.
Il fattore decisivo non è stata la velocità grezza di generazione del codice, ma la
combinazione di tre elementi: il \emph{workflow} \emph{Plan $\rightarrow$ Confirm $\rightarrow$ Implement}
(che ha evitato il \emph{sprawl}), le \emph{skill} dedicate (che hanno uniformato le
convenzioni tra moduli) e il \gls{mcp} (che ha abbattuto le allucinazioni su metriche e
dashboard).
